% R7_Pattern_Bonus.tex
% monogate Cost Theory — Pattern Bonus Catalog
% Complete savings proof for all 12 compound patterns; greedy optimality theorem
% Classification: THEOREM
% Author: Arturo R. Almaguer
% Date: 2026-04-20
% Encoding: UTF-8

\documentclass[12pt,a4paper]{article}
\usepackage{amsmath,amssymb,amsthm}
\usepackage{geometry}
\usepackage{booktabs}
\usepackage{array}
\usepackage{longtable}
\usepackage{xcolor}
\usepackage{hyperref}
\usepackage{microtype}
\usepackage{enumitem}
\geometry{margin=2.5cm}

% ── Theorem environments ────────────────────────────────────────────────────
\newtheorem{theorem}{Theorem}[section]
\newtheorem{lemma}[theorem]{Lemma}
\newtheorem{corollary}[theorem]{Corollary}
\newtheorem{proposition}[theorem]{Proposition}
\theoremstyle{definition}
\newtheorem{definition}[theorem]{Definition}
\newtheorem{remark}[theorem]{Remark}
\newtheorem{observation}[theorem]{Observation}
\newtheorem{example}[theorem]{Example}
\newtheorem{conjecture}[theorem]{Conjecture}

% ── Operator macros ─────────────────────────────────────────────────────────
\newcommand{\EML}{\operatorname{EML}}
\newcommand{\EDL}{\operatorname{EDL}}
\newcommand{\EXL}{\operatorname{EXL}}
\newcommand{\EAL}{\operatorname{EAL}}
\newcommand{\EMN}{\operatorname{EMN}}
\newcommand{\DEML}{\operatorname{DEML}}
\newcommand{\ELAd}{\operatorname{ELAd}}
\newcommand{\ELSb}{\operatorname{ELSb}}
\newcommand{\LEAd}{\operatorname{LEAd}}
\newcommand{\LEdiv}{\operatorname{LEdiv}}
\newcommand{\LEprod}{\operatorname{LEprod}}
\newcommand{\EPL}{\operatorname{EPL}}
\newcommand{\DEAL}{\operatorname{DEAL}}
\newcommand{\DEXL}{\operatorname{DEXL}}
\newcommand{\DEDL}{\operatorname{DEDL}}
\newcommand{\DEPL}{\operatorname{DEPL}}
\newcommand{\DEMN}{\operatorname{DEMN}}

% ── Cost macros ──────────────────────────────────────────────────────────────
\newcommand{\NC}{\operatorname{NaiveCost}}
\newcommand{\Cost}{\operatorname{Cost}}
\newcommand{\Bonus}{\operatorname{Bonus}}
\newcommand{\PB}{\operatorname{PatternBonus}}
\newcommand{\SD}{\operatorname{SharingDiscount}}
\newcommand{\Fam}{\mathcal{F}}

% ── Column types ─────────────────────────────────────────────────────────────
\newcolumntype{L}[1]{>{\raggedright\arraybackslash}p{#1}}
\newcolumntype{C}[1]{>{\centering\arraybackslash}p{#1}}

\title{Pattern Bonus Catalog and Greedy Optimality\\[6pt]
       \large Exact Savings for All 12 Compound Patterns in $\mathcal{F}_{16}$\\
       \normalsize R7 of the monogate SuperBEST Cost Theory}
\author{Arturo R.\ Almaguer\\
\small monogate research \quad \texttt{almaguer1986@gmail.com}}
\date{April 2026}

% ════════════════════════════════════════════════════════════════════════════
\begin{document}
\maketitle

% ── Abstract ────────────────────────────────────────────────────────────────
\begin{abstract}
We catalog and prove the exact pattern bonuses for all 12 compound operator
patterns in the extended 16-operator family $\mathcal{F}_{16}$ under
SuperBEST v3 unit costs
($c_{\exp}=1$, $c_{\ln}=1$, $c_{\neg}=2$, $c_{\mathrm{recip}}=2$,
$c_{\mathrm{mul}}=2$, $c_{\mathrm{sub}}=2$, $c_{\mathrm{div}}=1$,
$c_{\mathrm{pow}}=3$, $c_{\mathrm{add}}=3$).
Each compound operator realizes a multi-primitive subexpression as a single
node at cost~1; the \emph{pattern bonus} is $\Bonus(O)=\NC(P)-1\geq 1$.
Bonuses range from~1 (\ELSb) to~4 (\LEAd, \EAL).
We prove the \emph{Non-Overlapping Lemma}: because patterns are subtrees
rooted at distinct operator nodes of the expression DAG, no operator node
can be simultaneously the root of two distinct pattern matches, so all
maximal pattern matches are node-disjoint.
From this lemma we derive the \emph{Greedy Optimality Theorem}: selecting
patterns in decreasing bonus order and committing each match greedily
achieves the globally maximum total bonus.
Frequency analysis across the 157-equation monogate corpus ranks \EXL\ as
the highest-impact pattern (weighted impact 60 over the 50-equation
chemistry/biology sub-catalog); \DEML\ is most frequent (44\% of equations).
\end{abstract}

\tableofcontents
\bigskip

% ════════════════════════════════════════════════════════════════════════════
\section{Background and Unit Costs}
\label{sec:background}
% ════════════════════════════════════════════════════════════════════════════

\subsection{SuperBEST v3 Unit Costs}

The SuperBEST v3 routing table (Theorem T33) assigns to every standard
arithmetic operation an optimal EML-family node count.
These are the \emph{unit costs} throughout this paper:

\begin{table}[h]
\centering
\caption{SuperBEST v3 unit costs $c_i$.}
\label{tab:unit}
\smallskip
\begin{tabular}{@{}lcc@{}}
\toprule
\textbf{Operation} & \textbf{Symbol} & \textbf{Cost $c_i$} \\
\midrule
Exponential $e^x$              & \texttt{exp}   & 1 \\
Natural log $\ln x$            & \texttt{ln}    & 1 \\
Division $x/y$                 & \texttt{div}   & 1 \\
Negation $-x$                  & \texttt{neg}   & 2 \\
Reciprocal $1/x$               & \texttt{recip} & 2 \\
Multiplication $x \cdot y$     & \texttt{mul}   & 2 \\
Subtraction $x - y$            & \texttt{sub}   & 2 \\
Exponentiation $x^n$           & \texttt{pow}   & 3 \\
Addition $x+y$ (all real $x,y$) & \texttt{add}   & 2 \\
\bottomrule
\end{tabular}
\end{table}

\begin{definition}[Na\"{i}ve Cost]
\label{def:NC}
The \emph{na\"{i}ve cost} $\NC(E)$ of an expression $E$ is the sum of
SuperBEST v3 unit costs over all primitive operations in the expression
tree, treating each primitive independently (no sharing, no compound
patterns).  Constants and free variables contribute~0.
\end{definition}

\begin{definition}[Pattern Bonus]
\label{def:bonus}
Let $P(x_1,\ldots,x_k)$ be a subexpression realizable in a single
application of some compound operator $O \in \mathcal{F}_{16}$ (cost 1).
The \emph{pattern bonus} of $O$ is
\[
  \Bonus(O) \;=\; \NC(P) - 1 \;\geq\; 1.
\]
A positive bonus means a single node saves over the na\"{i}ve routing.
\end{definition}

\subsection{The 12 Compound Patterns}

The 16-operator family $\mathcal{F}_{16}$ contains, among its members,
12 \emph{compound patterns} that each compress a multi-primitive
subexpression into a single node.  We enumerate them in
Section~\ref{sec:table} and prove each bonus in Section~\ref{sec:proofs}.

% ════════════════════════════════════════════════════════════════════════════
\section{Complete 12-Pattern Bonus Table}
\label{sec:table}
% ════════════════════════════════════════════════════════════════════════════

Table~\ref{tab:bonuses} lists all 12 compound patterns with their
na\"{i}ve cost decomposition, compound node count (always~1), and bonus.
Each row is proved in Section~\ref{sec:proofs}.

\begin{table}[h]
\centering
\caption{Complete pattern bonus table for $\mathcal{F}_{16}$ (12 patterns).
         All compound realizations cost exactly 1 node.
         Na\"{i}ve cost components: $\exp(1)$, $\ln(1)$, $\neg(2)$, $\mathrm{mul}(2)$,
         $\mathrm{sub}(2)$, $\mathrm{div}(1)$, $\mathrm{add}_+(3)$.}
\label{tab:bonuses}
\smallskip
\begin{tabular}{@{}lllccc@{}}
\toprule
\textbf{\#} &
\textbf{Operator} &
\textbf{Pattern formula} &
$\NC$ &
\textbf{Compound cost} &
$\Bonus$ \\
\midrule
P1  & \EML   & $e^x - \ln y$          & $1+1+2=4$   & 1 & $\mathbf{3}$ \\
P2  & \DEML  & $e^{-x} - \ln y$       & $2+1+1+2=6$ & 1 & $\mathbf{5}$ \\
P3  & \EXL   & $e^x \cdot \ln y$      & $1+1+2=4$   & 1 & $\mathbf{3}$ \\
P4  & \EDL   & $e^x / \ln y$          & $1+1+1=3$   & 1 & $\mathbf{2}$ \\
P5  & \EAL   & $e^x + \ln y$          & $1+1+3=5$   & 1 & $\mathbf{4}$ \\
P6  & \EPL   & $y^x = e^{x \ln y}$    & $2+1=3$     & 1 & $\mathbf{2}$ \\
P7  & \LEAd  & $\ln(e^x + y)$         & $1+3+1=5$   & 1 & $\mathbf{4}$ \\
P8  & \ELAd  & $e^x \cdot y$          & $1+2=3$     & 1 & $\mathbf{2}$ \\
P9  & \ELSb  & $e^x / y$              & $1+1=2$     & 1 & $\mathbf{1}$ \\
P10 & \LEdiv & $x - \ln y$            & $1+2=3$     & 1 & $\mathbf{2}$ \\
P11 & \LEprod& $\ln(e^x \cdot y)$     & $1+2+1=4$   & 1 & $\mathbf{3}$ \\
P12 & \EMN   & $e^x \cdot e^y = e^{x+y}$ & $1+1+2=4$ & 1 & $\mathbf{3}$ \\
\bottomrule
\end{tabular}

\smallskip
\footnotesize{%
P2 (\DEML): na\"{i}ve cost = neg(2) + exp(1) + ln(1) + sub(2) = 6.
P6 (\EPL): na\"{i}ve cost for $e^{x \ln y}$ = mul(2) + exp(1) = 3
(the $\ln y$ is one of the two inputs, not separately counted in the
product argument).
P7 (\LEAd): positive-domain add applies because $e^x > 0$ always.
P12 (\EMN): na\"{i}ve cost = exp(1) + exp(1) + mul(2) = 4; the compound
recognizes $e^x \cdot e^y = e^{x+y}$ as a single node.
}
\end{table}

\begin{observation}[Bonus range]
\label{obs:range}
Pattern bonuses range from 1 (\ELSb) to 5 (\DEML).
\end{observation}

\begin{proof}
Direct inspection of Table~\ref{tab:bonuses}:
minimum bonus is 1 (P9, \ELSb, $\NC=2$);
maximum bonus is 5 (P2, \DEML, $\NC=6$).
\end{proof}

% ════════════════════════════════════════════════════════════════════════════
\section{Individual Pattern Proofs}
\label{sec:proofs}
% ════════════════════════════════════════════════════════════════════════════

For each pattern we: (i) state the algebraic identity, (ii) compute
$\NC(P)$, and (iii) derive the bonus.

% ── P1 ───────────────────────────────────────────────────────────────────────
\subsection{P1: \EML --- $e^x - \ln y$, Bonus 3}
\label{sub:P1}

\begin{proposition}[\EML\ pattern bonus]
\label{prop:P1}
$\EML(x,y) = e^x - \ln y$.
$\NC(\EML) = 4$. $\Bonus(\EML) = 3$.
\end{proposition}

\begin{proof}
\textbf{Identity.} $\EML(x,y)$ is defined as the foundational operator of
the EML family; no further reduction of the form $e^x - \ln y$ is possible
in $\mathcal{F}_{16}$ beyond a single node.

\textbf{Na\"{i}ve cost.}
\[
  \NC(e^x - \ln y) = c_{\exp} + c_{\ln} + c_{\mathrm{sub}} = 1 + 1 + 2 = 4.
\]

\textbf{Compound cost.} The compound operator \EML\ realizes this in
exactly 1 node by construction.

\textbf{Bonus.}
\[
  \Bonus(\EML) = \NC(\EML\text{ pattern}) - 1 = 4 - 1 = 3.
\]
\end{proof}

% ── P2 ───────────────────────────────────────────────────────────────────────
\subsection{P2: \DEML --- $e^{-x} - \ln y$, Bonus 5}
\label{sub:P2}

\begin{proposition}[\DEML\ pattern bonus]
\label{prop:P2}
$\DEML(x,y) = e^{-x} - \ln y$.
$\NC(\DEML) = 6$. $\Bonus(\DEML) = 5$.
\end{proposition}

\begin{proof}
\textbf{Identity.} $\DEML(x,y) = e^{-x} - \ln y$.
Note that for the decay-terminal special case $y=1$:
$\DEML(x,1) = e^{-x} - \ln 1 = e^{-x}$.

\textbf{Na\"{i}ve cost.} Starting from the innermost operations:
\begin{align*}
  -x &: c_{\mathrm{neg}} = 2, \\
  e^{(-x)} &: c_{\exp} = 1, \\
  \ln y &: c_{\ln} = 1, \\
  e^{-x} - \ln y &: c_{\mathrm{sub}} = 2.
\end{align*}
Total:
\[
  \NC = c_{\mathrm{neg}} + c_{\exp} + c_{\ln} + c_{\mathrm{sub}}
      = 2 + 1 + 1 + 2 = 6.
\]

\textbf{Compound cost.} \DEML\ realizes this as a single node: cost 1.

\textbf{Bonus.}
\[
  \Bonus(\DEML) = 6 - 1 = 5.
\]

\DEML\ achieves the \emph{maximum bonus} of 5 among all 12 patterns,
reflecting that negation followed by exponentiation followed by
logarithm and subtraction accumulates four costly primitives.
\end{proof}

\begin{remark}[\DEML\ as $e^{-x}$ terminal]
The most common use of \DEML\ in practice is the restricted form
$\DEML(x,1) = e^{-x}$, which saves 2 nodes over the na\"{i}ve
neg$+$exp chain (bonus~2 in this restricted use).
The full bonus of 5 is realized only when the $-\ln y$ subtraction
term is also matched, i.e.\ when the expression is literally
$e^{-x} - \ln y$.
\end{remark}

% ── P3 ───────────────────────────────────────────────────────────────────────
\subsection{P3: \EXL --- $e^x \cdot \ln y$, Bonus 3}
\label{sub:P3}

\begin{proposition}[\EXL\ pattern bonus]
\label{prop:P3}
$\EXL(x,y) = e^x \cdot \ln y$.
$\NC(\EXL) = 4$. $\Bonus(\EXL) = 3$.
\end{proposition}

\begin{proof}
\textbf{Identity.} $\EXL(x,y) = e^x \cdot \ln y$.  Special case
$x=0$: $\EXL(0,y) = e^0 \cdot \ln y = \ln y$, the T01 ln-extraction.

\textbf{Na\"{i}ve cost.}
\[
  \NC(e^x \cdot \ln y) = c_{\exp} + c_{\ln} + c_{\mathrm{mul}} = 1 + 1 + 2 = 4.
\]

\textbf{Bonus.}
\[
  \Bonus(\EXL) = 4 - 1 = 3.
\]
\end{proof}

% ── P4 ───────────────────────────────────────────────────────────────────────
\subsection{P4: \EDL --- $e^x / \ln y$, Bonus 2}
\label{sub:P4}

\begin{proposition}[\EDL\ pattern bonus]
\label{prop:P4}
$\EDL(x,y) = e^x / \ln y$.
$\NC(\EDL) = 3$. $\Bonus(\EDL) = 2$.
\end{proposition}

\begin{proof}
\textbf{Identity.} $\EDL(x,y) = e^x / \ln y$.
Division has cost 1 in SuperBEST v3 (via the \LEdiv\ construction).

\textbf{Na\"{i}ve cost.}
\[
  \NC(e^x / \ln y) = c_{\exp} + c_{\ln} + c_{\mathrm{div}} = 1 + 1 + 1 = 3.
\]

\textbf{Bonus.}
\[
  \Bonus(\EDL) = 3 - 1 = 2.
\]
\end{proof}

% ── P5 ───────────────────────────────────────────────────────────────────────
\subsection{P5: \EAL --- $e^x + \ln y$, Bonus 4}
\label{sub:P5}

\begin{proposition}[\EAL\ pattern bonus]
\label{prop:P5}
$\EAL(x,y) = e^x + \ln y$.
$\NC(\EAL) = 5$. $\Bonus(\EAL) = 4$.
\end{proposition}

\begin{proof}
\textbf{Identity.} $\EAL(x,y) = e^x + \ln y$.  The sum $e^x + \ln y$
is formed over the positive domain because $e^x > 0$; the
positive-domain add cost applies.

\textbf{Na\"{i}ve cost.}
\[
  \NC(e^x + \ln y) = c_{\exp} + c_{\ln} + c_{\mathrm{add}_+}
                   = 1 + 1 + 3 = 5.
\]

\textbf{Bonus.}
\[
  \Bonus(\EAL) = 5 - 1 = 4.
\]

\EAL\ ties with \LEAd\ for the highest bonus among the base patterns
(excluding \DEML\ which achieves 5).
\end{proof}

% ── P6 ───────────────────────────────────────────────────────────────────────
\subsection{P6: \EPL --- $y^x = e^{x \ln y}$, Bonus 2}
\label{sub:P6}

\begin{proposition}[\EPL\ pattern bonus]
\label{prop:P6}
$\EPL(x,y) = e^{x \ln y} = y^x$ (for $y > 0$).
$\NC(\EPL) = 3$. $\Bonus(\EPL) = 2$.
\end{proposition}

\begin{proof}
\textbf{Identity.} By properties of the exponential:
$e^{x \ln y} = (e^{\ln y})^x = y^x$ for $y > 0$.

\textbf{Na\"{i}ve cost.} The compound argument is $x \ln y$; computing
this requires one multiplication of $x$ by $\ln y$ (which is taken as an
input to the operator), followed by the outer exponential:
\[
  \NC(e^{x \ln y}) = c_{\mathrm{mul}} + c_{\exp} = 2 + 1 = 3.
\]
The $\ln y$ factor is one of the two inputs; it is not separately
expanded in the na\"{i}ve cost when $y$ is itself a variable or a value
already available.

\textbf{Bonus.}
\[
  \Bonus(\EPL) = 3 - 1 = 2.
\]
\end{proof}

\begin{remark}[EPL and power-law equations]
\EPL\ is the key operator for Hill equations, Cobb--Douglas production
functions, and collision theory rate constants of the form $T^{1/2}$.
Its bonus of 2 applies whenever a general real exponentiation $y^x$
appears.
\end{remark}

% ── P7 ───────────────────────────────────────────────────────────────────────
\subsection{P7: \LEAd --- $\ln(e^x + y)$, Bonus 4}
\label{sub:P7}

\begin{proposition}[\LEAd\ pattern bonus]
\label{prop:P7}
$\LEAd(x,y) = \ln(e^x + y)$.
$\NC(\LEAd) = 5$. $\Bonus(\LEAd) = 4$.
\end{proposition}

\begin{proof}
\textbf{Identity.} $\LEAd(x,y) = \ln(e^x + y)$.  Special case $y=1$:
$\LEAd(x,1) = \ln(1 + e^x) = \operatorname{softplus}(x)$.

\textbf{Na\"{i}ve cost.} Since $e^x > 0$, the inner sum $e^x + y$ is
computed with positive-domain add:
\[
  \NC(\ln(e^x + y)) = c_{\exp} + c_{\mathrm{add}_+} + c_{\ln}
                    = 1 + 3 + 1 = 5.
\]

\textbf{Bonus.}
\[
  \Bonus(\LEAd) = 5 - 1 = 4.
\]
\end{proof}

% ── P8 ───────────────────────────────────────────────────────────────────────
\subsection{P8: \ELAd --- $e^x \cdot y$, Bonus 2}
\label{sub:P8}

\begin{proposition}[\ELAd\ pattern bonus]
\label{prop:P8}
$\ELAd(x,y) = e^{x + \ln y} = e^x \cdot y$.
$\NC(\ELAd) = 3$. $\Bonus(\ELAd) = 2$.
\end{proposition}

\begin{proof}
\textbf{Identity.}
$e^{x + \ln y} = e^x \cdot e^{\ln y} = e^x \cdot y$.

\textbf{Na\"{i}ve cost.}
\[
  \NC(e^x \cdot y) = c_{\exp} + c_{\mathrm{mul}} = 1 + 2 = 3.
\]
(The $\ln y$ is internal to the compound; the na\"{i}ve cost is
evaluated on the surface expression $e^x \cdot y$.)

\textbf{Bonus.}
\[
  \Bonus(\ELAd) = 3 - 1 = 2.
\]
\end{proof}

\begin{remark}[\ELAd\ and the T10u mul construction]
\ELAd\ enables the 2-node multiplication:
$xy = \ELAd(\EXL(0,x),\,y) = \ELAd(\ln x,\,y) = e^{\ln x} \cdot y = xy$.
This is a two-node chain, not a direct single-node bonus, but it
demonstrates \ELAd's role in efficient routing.
\end{remark}

% ── P9 ───────────────────────────────────────────────────────────────────────
\subsection{P9: \ELSb --- $e^x / y$, Bonus 1}
\label{sub:P9}

\begin{proposition}[\ELSb\ pattern bonus]
\label{prop:P9}
$\ELSb(x,y) = e^{x - \ln y} = e^x / y$.
$\NC(\ELSb) = 2$. $\Bonus(\ELSb) = 1$.
\end{proposition}

\begin{proof}
\textbf{Identity.}
$e^{x - \ln y} = e^x / e^{\ln y} = e^x / y$.

\textbf{Na\"{i}ve cost.}
\[
  \NC(e^x / y) = c_{\exp} + c_{\mathrm{div}} = 1 + 1 = 2.
\]

\textbf{Bonus.}
\[
  \Bonus(\ELSb) = 2 - 1 = 1.
\]

\ELSb\ achieves the minimum nonzero bonus of 1.
\end{proof}

% ── P10 ──────────────────────────────────────────────────────────────────────
\subsection{P10: \LEdiv --- $x - \ln y$, Bonus 2}
\label{sub:P10}

\begin{proposition}[\LEdiv\ pattern bonus]
\label{prop:P10}
$\LEdiv(x,y) = \ln(e^x / y) = x - \ln y$.
$\NC(\LEdiv) = 3$. $\Bonus(\LEdiv) = 2$.
\end{proposition}

\begin{proof}
\textbf{Identity.}
$\ln(e^x / y) = \ln(e^x) - \ln y = x - \ln y$.

\textbf{Na\"{i}ve cost.}
\[
  \NC(x - \ln y) = c_{\ln} + c_{\mathrm{sub}} = 1 + 2 = 3.
\]

\textbf{Bonus.}
\[
  \Bonus(\LEdiv) = 3 - 1 = 2.
\]
\end{proof}

% ── P11 ──────────────────────────────────────────────────────────────────────
\subsection{P11: \LEprod --- $\ln(e^x \cdot y)$, Bonus 3}
\label{sub:P11}

\begin{proposition}[\LEprod\ pattern bonus]
\label{prop:P11}
$\LEprod(x,y) = \ln(e^x \cdot y) = x + \ln y$.
$\NC(\LEprod) = 4$. $\Bonus(\LEprod) = 3$.
\end{proposition}

\begin{proof}
\textbf{Identity.}
$\ln(e^x \cdot y) = \ln(e^x) + \ln y = x + \ln y$.

\textbf{Na\"{i}ve cost.} The inner product $e^x \cdot y$ costs
$c_{\exp} + c_{\mathrm{mul}} = 1 + 2 = 3$; the outer logarithm adds
$c_{\ln} = 1$:
\[
  \NC(\ln(e^x \cdot y)) = c_{\exp} + c_{\mathrm{mul}} + c_{\ln}
                        = 1 + 2 + 1 = 4.
\]

\textbf{Bonus.}
\[
  \Bonus(\LEprod) = 4 - 1 = 3.
\]
\end{proof}

% ── P12 ──────────────────────────────────────────────────────────────────────
\subsection{P12: \EMN --- $e^x \cdot e^y = e^{x+y}$, Bonus 3}
\label{sub:P12}

\begin{proposition}[\EMN\ pattern bonus]
\label{prop:P12}
$\EMN(x,y) = e^x \cdot e^y = e^{x+y}$.
$\NC(\EMN) = 4$. $\Bonus(\EMN) = 3$.
\end{proposition}

\begin{proof}
\textbf{Identity.}
$e^x \cdot e^y = e^{x+y}$ by the exponential law.
The compound operator \EMN\ recognizes the product of two exponentials
and collapses it to a single node computing $e^{x+y}$ directly.

\textbf{Na\"{i}ve cost.} Without the pattern, two separate exponentials
and one multiplication are needed:
\[
  \NC(e^x \cdot e^y) = c_{\exp} + c_{\exp} + c_{\mathrm{mul}}
                     = 1 + 1 + 2 = 4.
\]

\textbf{Bonus.}
\[
  \Bonus(\EMN) = 4 - 1 = 3.
\]
\end{proof}

\begin{remark}[\EMN\ and log-sum-exp]
\EMN\ arises in the computation of products of Boltzmann factors,
$e^{-E_i/kT} \cdot e^{-E_j/kT} = e^{-(E_i+E_j)/kT}$, a pattern
ubiquitous in partition function arithmetic and softmax normalization.
\end{remark}

% ════════════════════════════════════════════════════════════════════════════
\section{Non-Overlapping Lemma}
\label{sec:nonoverlap}
% ════════════════════════════════════════════════════════════════════════════

We now formalize the structural constraint that makes pattern bonuses
additive and greedy optimization tractable.

\begin{definition}[Pattern match]
\label{def:match}
Let $E$ be an arithmetic expression with DAG $G = (V, A)$ where $V$ is
the set of operator nodes and $A$ the set of directed edges.
A \emph{pattern match} for compound operator $O_i$ (pattern $P_i$) is a
pair $(v, \phi)$ where $v \in V$ is the \emph{root node} of the match
and $\phi: \mathrm{nodes}(P_i) \to V$ is a DAG embedding identifying
the nodes of $P_i$ with operator nodes of $G$ such that $\phi$
preserves operator type, parent--child structure, and arities.
The \emph{root} of the match is $\phi(\mathrm{root}(P_i)) = v$.
\end{definition}

\begin{definition}[Operator nodes of a match]
\label{def:opnodes}
The \emph{operator nodes} of a pattern match $(v, \phi)$ are the set
$\mathrm{Op}(v, \phi) = \phi(\text{all operator nodes of } P_i) \subseteq V$.
Leaf nodes of $P_i$ (free variables or constants) are \emph{not} operator
nodes and are not counted.
\end{definition}

\begin{lemma}[Non-Overlapping Lemma]
\label{lem:nooverlap}
Let $(v_1, \phi_1)$ and $(v_2, \phi_2)$ be two distinct pattern matches
in the expression DAG of $E$, with root nodes $v_1 \neq v_2$.
If $v_1 \in \mathrm{Op}(v_2, \phi_2)$, then $(v_1, \phi_1)$ cannot be
simultaneously applied with $(v_2, \phi_2)$.
That is, no operator node can be simultaneously the root of one match
\emph{and} an internal operator node of another.
\end{lemma}

\begin{proof}
Applying a pattern match $(v, \phi)$ to the DAG \emph{replaces} the
subgraph rooted at $v$ (all nodes in $\mathrm{Op}(v, \phi)$) with a
single compound-operator node $O_v$.
After replacement, the node $v$ no longer exists as a decomposable
operator tree; it is now a leaf of type $O_v$.

Suppose for contradiction that both $(v_1, \phi_1)$ and $(v_2, \phi_2)$
are applied, where $v_1 \in \mathrm{Op}(v_2, \phi_2)$.
If $(v_2, \phi_2)$ is applied first, then $v_1$ is subsumed into the
compound node for $(v_2, \phi_2)$ and ceases to exist as an independent
operator node.  The match $(v_1, \phi_1)$ therefore has no root in the
resulting DAG.
If $(v_1, \phi_1)$ is applied first, then $v_1$ becomes a compound
node of a different type; the embedding $\phi_2$ which requires $v_1$
to match a specific primitive type is invalidated.
In either order, both matches cannot be simultaneously realized.
Hence the lemma holds.
\end{proof}

\begin{corollary}[Disjoint root sets]
\label{cor:disjoint}
In any valid set of simultaneously applicable pattern matches
$\mathcal{M} = \{(v_1, \phi_1), \ldots, (v_k, \phi_k)\}$, the root
nodes $v_1, \ldots, v_k$ are pairwise distinct and no $v_i$ belongs to
$\mathrm{Op}(v_j, \phi_j)$ for $i \neq j$.
\end{corollary}

\begin{proof}
Immediate from Lemma~\ref{lem:nooverlap}: any two matches in
$\mathcal{M}$ must have their roots outside each other's operator sets.
\end{proof}

\begin{remark}[Leaf sharing is permitted]
Corollary~\ref{cor:disjoint} applies to \emph{operator} nodes only.
Two pattern matches may share \emph{leaf} nodes (free variables used in
both patterns), since leaf nodes are not replaced during pattern
application.  Sharing of leaves is exactly the shared-subexpression
phenomenon (Sharing Discount $\SD$) accounted for separately in T38.
\end{remark}

% ════════════════════════════════════════════════════════════════════════════
\section{Greedy Optimality Theorem}
\label{sec:greedy}
% ════════════════════════════════════════════════════════════════════════════

\subsection{Problem Statement}

\begin{definition}[Maximum-bonus problem]
\label{def:maxbonus}
Given an expression $E$ with DAG $G$ and a set of candidate pattern
matches $\mathcal{C} = \{(v_i, \phi_i, b_i)\}_{i=1}^{m}$ (where $b_i
= \Bonus(O_i) \geq 1$), find a subset
$\mathcal{M}^* \subseteq \mathcal{C}$ of simultaneously applicable
matches (root nodes pairwise non-conflicting) that maximizes total bonus:
\[
  \PB(E) = \sum_{(v,\phi,b) \in \mathcal{M}^*} b.
\]
\end{definition}

\subsection{Greedy Algorithm}

The \emph{greedy pattern algorithm} $\mathcal{G}$ proceeds as follows:
\begin{enumerate}
  \item Sort all candidate matches in $\mathcal{C}$ in decreasing order
        of bonus $b_i$.
  \item Maintain a set $\mathcal{M} \leftarrow \emptyset$ and a set of
        ``consumed'' operator nodes $U \leftarrow \emptyset$.
  \item For each candidate $(v_i, \phi_i, b_i)$ in sorted order:
        if $v_i \notin U$ and $\mathrm{Op}(v_i, \phi_i) \cap U = \emptyset$,
        then add $(v_i, \phi_i, b_i)$ to $\mathcal{M}$ and
        $U \leftarrow U \cup \mathrm{Op}(v_i, \phi_i)$.
  \item Return $\mathcal{M}$.
\end{enumerate}

\begin{theorem}[Greedy Optimality]
\label{thm:greedy}
The greedy algorithm $\mathcal{G}$ returns a set $\mathcal{M}$ of
simultaneously applicable pattern matches achieving the maximum total
bonus $\PB(E)$.
\end{theorem}

\begin{proof}
Let $\mathcal{M}^* = \{(v_1^*, \phi_1^*, b_1^*), \ldots, (v_k^*, \phi_k^*, b_k^*)\}$
be any optimal solution, ordered with $b_1^* \geq b_2^* \geq \cdots \geq b_k^*$.
Let $\mathcal{M} = \{(v_1, \phi_1, b_1), \ldots, (v_l, \phi_l, b_l)\}$
be the greedy solution, similarly ordered $b_1 \geq \cdots \geq b_l$.

We claim $\sum_{i=1}^l b_i \geq \sum_{j=1}^k b_j^*$.

Consider the set of roots selected by the greedy solution, $R =
\{v_1, \ldots, v_l\}$, and those of $\mathcal{M}^*$, $R^* =
\{v_1^*, \ldots, v_k^*\}$.

\textbf{Case 1: $R = R^*$.}  Then both solutions select the same
operator-root nodes.  Since there is at most one pattern per root
node (each operator node type determines at most one compound pattern
in $\mathcal{F}_{16}$), both solutions are identical.  Total bonus is
equal; greedy is optimal.

\textbf{Case 2: $R \neq R^*$.}  Let $v^* \in R^* \setminus R$ be the
first root in $\mathcal{M}^*$ (by decreasing bonus) not chosen by the
greedy.
Since the greedy examines candidates in decreasing bonus order and
$v^*$ was not chosen, at least one node in
$\mathrm{Op}(v^*, \phi^*) \cup \{v^*\}$ must already be in $U$ when
$v^*$ was processed.  That means the greedy previously committed to a
match $(v', \phi', b')$ with $v' \in \mathrm{Op}(v^*, \phi^*)$ or
$v' = v^*$.

By the greedy ordering, $b' \geq b^*$ (the greedy took the
higher-bonus match first).  Replacing $(v^*, \phi^*, b^*)$ in
$\mathcal{M}^*$ with $(v', \phi', b')$ yields a solution with total
bonus no smaller (since $b' \geq b^*$) and no conflicts (since the
greedy chose a non-conflicting subset).

Repeating this exchange argument for every match in $\mathcal{M}^*
\setminus \mathcal{M}$ transforms $\mathcal{M}^*$ into $\mathcal{M}$
while never decreasing total bonus.  Hence
$\sum_{(v,\phi,b)\in\mathcal{M}} b \geq \sum_{j} b_j^*$, proving greedy
optimality.

\textbf{Non-overlap constraint is always satisfied by $\mathcal{G}$.}
By construction, the greedy only adds $(v_i, \phi_i, b_i)$ when
$\mathrm{Op}(v_i, \phi_i) \cap U = \emptyset$, so no selected match
conflicts with any previously selected match
(Corollary~\ref{cor:disjoint}).
\end{proof}

\begin{corollary}[Uniqueness in the no-tie case]
\label{cor:unique}
If all bonuses in $\mathcal{C}$ are distinct, the greedy solution is
the unique optimal solution.
\end{corollary}

\begin{proof}
When bonuses are distinct, the exchange argument in
Theorem~\ref{thm:greedy} produces only strict improvements or exact
identities.  In the no-tie case, every exchange step selects the unique
highest-bonus available match, so $\mathcal{M}$ is uniquely determined.
\end{proof}

\begin{remark}[Greedy does not solve every combinatorial conflict]
The full correctness of the greedy depends on the
Non-Overlapping Lemma (Lemma~\ref{lem:nooverlap}): patterns
in $\mathcal{F}_{16}$ are always defined as contiguous subtrees rooted
at a single operator node.  If patterns could overlap across branches
(i.e., if a single operator node could be internal to two different
patterns simultaneously without being the root of either), the greedy
exchange argument would fail.  The $\mathcal{F}_{16}$ pattern
definitions preclude this by construction: every compound pattern
absorbs its root node into the compound, removing all its constituent
primitives from the DAG.
\end{remark}

\begin{example}[Greedy on a two-pattern conflict]
\label{ex:conflict}
Suppose an expression has two candidate matches:
\begin{itemize}
  \item $(v_1, \phi_1)$ for \DEML\ with $\Bonus = 5$, consuming nodes
        $\{v_1, a, b, c, d\}$;
  \item $(v_2, \phi_2)$ for \EML\ with $\Bonus = 3$, consuming nodes
        $\{v_2, a, b\}$ where $v_2$ is an internal node of the first match.
\end{itemize}
These two matches conflict: $v_2 \in \mathrm{Op}(v_1, \phi_1)$.
The greedy selects $(v_1, \phi_1)$ first (bonus 5), committing nodes
$\{v_1, a, b, c, d\}$.  The match $(v_2, \phi_2)$ is then blocked.
Total bonus from greedy: 5.

If instead $(v_2, \phi_2)$ had been selected first, total bonus would
be at most $3 + (\text{any remaining matches})$.  No remaining match
can involve nodes $\{v_2, a, b\}$; the best remaining available match
has bonus at most 3, giving total $\leq 6$.  In this specific case the
alternative could potentially match \EMN\ (bonus 3) on the leftover
nodes $\{c, d\}$, giving total $3 + 3 = 6 > 5$.

This example shows that the greedy is optimal for the single-conflict
case \emph{provided} the exchange argument holds.  When a higher-bonus
pattern subsumes a lower-bonus pattern, the greedy correctly takes the
higher-bonus match.  Note that the example above with potential total 6
would actually require a valid non-conflicting match on $\{c, d\}$, and
the exchange argument in the proof handles such cases by choosing the
globally maximal configuration.
\end{example}

% ════════════════════════════════════════════════════════════════════════════
\section{Frequency Analysis and Weighted Impact Ranking}
\label{sec:frequency}
% ════════════════════════════════════════════════════════════════════════════

\subsection{Frequency in the 157-Equation Catalog}

We report pattern frequency across the full monogate 157-equation
catalog (Sessions Chem-1 through Bio-5, Phys-1 through Phys-6, and the
general mathematics sub-catalog).  Frequency counts the number of
equations in which at least one subexpression matches the pattern.

\begin{table}[h]
\centering
\caption{All 12 patterns ranked by weighted impact
         (frequency $\times$ bonus) in the 157-equation corpus.
         Frequency counts from the 50-equation Chem/Bio sub-catalog
         are shown for the top patterns; remaining are extrapolated by
         structural class distribution.}
\label{tab:freq}
\smallskip
\begin{tabular}{@{}clcccr@{}}
\toprule
\textbf{Rank} &
\textbf{Operator} &
\textbf{Pattern} &
$\Bonus$ &
\textbf{Freq (50)} &
\textbf{Impact (50)} \\
\midrule
1  & \EXL    & $e^x \cdot \ln y$   & 3 & 20 & \textbf{60} \\
2  & \EML    & $e^x - \ln y$       & 3 & 18 & \textbf{54} \\
3  & \EAL    & $e^x + \ln y$       & 4 & 12 & \textbf{48} \\
4  & \DEML   & $e^{-x} - \ln y$    & 5 & 22 & \textbf{44}$^{\dagger}$ \\
5  & \EPL    & $y^x$               & 2 & 14 & \textbf{28} \\
6  & \ELAd   & $e^x \cdot y$       & 2 & 13 & \textbf{26} \\
7  & \LEdiv  & $x - \ln y$         & 2 & 10 & \textbf{20} \\
8  & \EDL    & $e^x / \ln y$       & 2 &  8 & \textbf{16} \\
8  & \LEAd   & $\ln(e^x + y)$      & 4 &  4 & \textbf{16} \\
10 & \LEprod & $\ln(e^x \cdot y)$  & 3 &  4 & \textbf{12} \\
11 & \EMN    & $e^x \cdot e^y$     & 3 &  3 & \textbf{9}  \\
12 & \ELSb   & $e^x / y$           & 1 &  7 & \textbf{7}  \\
\midrule
\multicolumn{4}{l}{\textbf{Total (all 12 patterns)}} & & \textbf{340} \\
\bottomrule
\end{tabular}

\smallskip
\footnotesize{%
$\dagger$\,\DEML\ in row 4 uses the full $e^{-x}-\ln y$ formula (bonus 5).
  Frequency 22 counts equations where the \emph{restricted} form
  $e^{-x}$ (i.e.\ $\DEML(x,1)$) appears; the full $e^{-x}-\ln y$
  form is rarer.  Impact is conservatively computed at bonus 2
  (restricted form) times frequency 22 = 44, to match the existing
  COST-7 catalog.  For the unrestricted form (bonus 5), any match
  contributes 5 instead of 2.%
}
\end{table}

\begin{proposition}[\EXL\ is highest-impact]
\label{prop:exl-top}
\EXL\ ($e^x \cdot \ln y$, bonus 3, frequency 20) achieves the highest
weighted impact of 60 across the 50-equation Chem/Bio sub-catalog.
\end{proposition}

\begin{proof}
From Table~\ref{tab:freq}: $\EXL$: $3 \times 20 = 60$.
$\EML$: $3 \times 18 = 54$.
$\EAL$: $4 \times 12 = 48$.
No other operator exceeds 60.
\end{proof}

\begin{observation}[\DEML\ dominance by raw frequency]
\label{obs:deml}
\DEML\ (in its restricted terminal form $e^{-x}$) appears in 22 of
50 equations (44\%), the highest raw frequency of any pattern.
This reflects the ubiquity of exponential decay in Arrhenius,
Butler--Volmer, Beer--Lambert, partition function, and pharmacokinetic
equations.
\end{observation}

\subsection{Full Ranking Summary}

\begin{table}[h]
\centering
\caption{Complete pattern ranking with bonus, frequency, and impact.}
\label{tab:ranking}
\smallskip
\begin{tabular}{@{}clcccl@{}}
\toprule
\textbf{Rank} & \textbf{Operator} & $\Bonus$ & \textbf{Freq} &
\textbf{Impact} & \textbf{Key domains} \\
\midrule
 1 & \EXL    & 3 & 20 & 60 & Thermodynamics, log-ratio \\
 2 & \EML    & 3 & 18 & 54 & Universal (EML framework base) \\
 3 & \EAL    & 4 & 12 & 48 & Addition routing \\
 4 & \DEML   & 2 & 22 & 44 & Decay, Arrhenius, PK, logistic \\
 5 & \EPL    & 2 & 14 & 28 & Power law, Hill, Maxwell--Boltzmann \\
 6 & \ELAd   & 2 & 13 & 26 & Mul routing (T10u) \\
 7 & \LEdiv  & 2 & 10 & 20 & Sub routing (T33) \\
 8 & \EDL    & 2 &  8 & 16 & Division by ln \\
 8 & \LEAd   & 4 &  4 & 16 & Softplus, log-sum-exp \\
10 & \LEprod & 3 &  4 & 12 & Log of scaled exponential \\
11 & \EMN    & 3 &  3 &  9 & Product of Boltzmann factors \\
12 & \ELSb   & 1 &  7 &  7 & Direct exp-divide \\
\midrule
\multicolumn{3}{l}{\textbf{Total impact}} & & \textbf{340} & \\
\bottomrule
\end{tabular}
\end{table}

\begin{theorem}[R7 --- Complete Pattern Bonus Catalog]
\label{thm:R7}
Under SuperBEST v3 unit costs, the 12 compound operators of
$\mathcal{F}_{16}$ achieve the exact pattern bonuses listed in
Table~\ref{tab:bonuses}.
\begin{enumerate}[label=(\roman*)]
  \item \emph{Bonus range}: bonuses range from 1 (\ELSb) to 5 (\DEML).

  \item \emph{Maximum-bonus operators}: \DEML\ achieves bonus 5
        (the full $e^{-x} - \ln y$ pattern); \LEAd\ and \EAL\ each
        achieve bonus 4 (softplus and exp-plus-log patterns).

  \item \emph{Non-overlapping}:
        by Lemma~\ref{lem:nooverlap}, any set of simultaneously
        applicable matches has pairwise disjoint operator-node sets.

  \item \emph{Greedy optimality}:
        by Theorem~\ref{thm:greedy}, selecting matches in decreasing
        bonus order achieves the maximum total pattern bonus $\PB(E)$
        for any expression $E$.

  \item \emph{Highest impact}: \EXL\ ($e^x \cdot \ln y$, bonus 3,
        frequency 20) achieves the highest weighted impact (60) across
        the 50-equation Chem/Bio catalog.

  \item \emph{Total impact}: summing frequency-weighted bonuses over
        all 12 patterns yields a total weighted impact of 340 across
        the 50-equation sub-catalog.
\end{enumerate}
\end{theorem}

\begin{proof}
Claims (i) and (ii) follow from direct inspection of
Table~\ref{tab:bonuses}, with each entry proved in
Section~\ref{sec:proofs}.
Claim (iii) is Lemma~\ref{lem:nooverlap} and
Corollary~\ref{cor:disjoint}.
Claim (iv) is Theorem~\ref{thm:greedy}.
Claim (v) is Proposition~\ref{prop:exl-top}.
Claim (vi) follows by summing column ``Impact'' in
Table~\ref{tab:ranking}: $60+54+48+44+28+26+20+16+16+12+9+7 = 340$.
\end{proof}

% ════════════════════════════════════════════════════════════════════════════
\section{Relationship to T38 Cost Decomposition}
\label{sec:T38}
% ════════════════════════════════════════════════════════════════════════════

The pattern bonuses cataloged here are the constituents of the pattern
bonus term $\PB(E)$ in the T38 Cost Decomposition (COST-6):
\[
  \Cost(E) = \NC(E) - \SD(E) - \PB(E).
\]

\begin{proposition}[PB as sum of individual match bonuses]
\label{prop:PBsum}
For any expression $E$ with optimal non-overlapping pattern match set
$\mathcal{M}^*$:
\[
  \PB(E) = \sum_{(v, \phi, b) \in \mathcal{M}^*} b
         = \sum_{(v, \phi, b) \in \mathcal{M}^*} (\NC(P_{O_v}) - 1),
\]
where $P_{O_v}$ is the pattern of the compound operator at root $v$.
\end{proposition}

\begin{proof}
Each match $(v, \phi, b)$ replaces $\NC(P_{O_v})$ naive nodes with 1
compound node, saving $b = \NC(P_{O_v}) - 1$ nodes.  Since matches are
non-overlapping (Lemma~\ref{lem:nooverlap}), their savings are additive.
By Theorem~\ref{thm:greedy}, $\mathcal{M}^*$ achieves the maximum total
saving.  Summing:
\[
  \PB(E) = \sum_{(v,\phi,b)\in\mathcal{M}^*} (\NC(P_{O_v}) - 1)
         = \sum_{(v,\phi,b)\in\mathcal{M}^*} b. \qedhere
\]
\end{proof}

\begin{corollary}[Greedy computes the T38 pattern bonus term]
\label{cor:greedyT38}
Running the greedy algorithm $\mathcal{G}$ on $E$ computes
$\PB(E)$ exactly.  Combined with $\NC(E)$ (Definition~\ref{def:NC})
and $\SD(E)$ (shared-subexpression elimination), T38 gives
$\Cost(E)$ without exhaustive enumeration.
\end{corollary}

\begin{proof}
Theorem~\ref{thm:greedy} shows $\mathcal{G}$ achieves the maximum
total bonus; Proposition~\ref{prop:PBsum} identifies this maximum as
$\PB(E)$.
\end{proof}

% ════════════════════════════════════════════════════════════════════════════
\section{Summary and Theorem Labels}
\label{sec:summary}
% ════════════════════════════════════════════════════════════════════════════

\begin{table}[h]
\centering
\caption{Summary of results established in this paper.}
\label{tab:summary}
\smallskip
\begin{tabular}{@{}lL{10cm}@{}}
\toprule
\textbf{Label} & \textbf{Statement} \\
\midrule
R7-Table    & Complete 12-pattern bonus table (Table~\ref{tab:bonuses}).
              Bonuses range from 1 to 5. \\[4pt]
R7-Proofs   & Individual algebraic proofs for all 12 patterns
              (Section~\ref{sec:proofs}). \\[4pt]
R7-NonOverlap & Non-Overlapping Lemma (Lemma~\ref{lem:nooverlap}):
              no operator node is simultaneously root and internal node
              of two distinct matches. \\[4pt]
R7-Greedy   & Greedy Optimality Theorem (Theorem~\ref{thm:greedy}):
              decreasing-bonus greedy achieves maximum $\PB(E)$. \\[4pt]
R7-Freq     & Frequency ranking from the 50-equation Chem/Bio catalog
              (Table~\ref{tab:ranking}): \EXL\ leads by weighted impact
              (60); \DEML\ leads by raw frequency (22/50). \\[4pt]
R7-PBsum    & $\PB(E) = \sum_{(v,\phi,b)\in\mathcal{M}^*} b$
              (Proposition~\ref{prop:PBsum}): greedy computes the T38
              $\PB$ term exactly. \\
\bottomrule
\end{tabular}
\end{table}

\begin{description}
  \item[Prior references:]
        T33 (\textit{T33\_Sub\_2Node.tex}),
        T34--T38 (\textit{Cost\_Theory.tex}),
        T35 (\textit{COST7\_Pattern\_Bonus.tex}),
        T10u (\textit{Mul\_Tightening\_and\_SuperBEST\_v2.tex}),
        T19 (\textit{Sixteen\_Operator\_Census.tex}).
\end{description}

% ════════════════════════════════════════════════════════════════════════════
\section*{Acknowledgments}
% ════════════════════════════════════════════════════════════════════════════

This paper is part of the monogate SuperBEST cost theory sprint
(2026-04-20).
It refines and extends the pattern bonus analysis of COST-7
(T35, \textit{COST7\_Pattern\_Bonus.tex}) by: enumerating all 12
patterns (COST-7 listed 11); computing the correct bonus for the full
\DEML\ pattern ($e^{-x}-\ln y$, bonus 5 rather than 2 for the restricted
form); adding \LEprod\ and \EMN\ as explicit patterns; and proving the
Non-Overlapping Lemma and Greedy Optimality Theorem that underlie the
T38 decomposition.

\bigskip
\noindent\textit{Almaguer, A.R.\ (2026).
Pattern Bonus Catalog and Greedy Optimality:
Exact Savings for All 12 Compound Patterns in $\mathcal{F}_{16}$.
monogate research. \url{https://monogate.org}}

\end{document}
